\documentclass[11pt]{article}
\usepackage[T1]{fontenc}
\usepackage{lmodern,microtype,amsmath,amssymb}
\usepackage[a4paper,margin=28mm,top=24mm,bottom=24mm]{geometry}
\usepackage[hidelinks]{hyperref}
\hypersetup{pdftitle={Critical points and sphere insertions},
 pdfsubject={Closed amplitudes and marked reconstruction in Landau--Ginzburg theory}}
\newcommand{\Tr}{\operatorname{Tr}}
\newcommand{\Res}{\operatorname{Res}}
\setlength{\parindent}{1.1em}
\setlength{\parskip}{0.18em}
\raggedbottom
\begin{document}
\begin{center}
{\LARGE Critical points and sphere insertions\par}
\smallskip
{\large Closed amplitudes in a Landau--Ginzburg family\par}
\end{center}

\noindent
Consider the cubic potential \(W_t=x^3/3-tx\). For \(t\ne0\) its two
critical points are \(\sqrt t\) and \(-\sqrt t\). In the topologically
twisted theory the sphere weights of these vacua are
\(1/(2\sqrt t)\) and \(-1/(2\sqrt t)\). Each diverges as the vacua
meet. Their sum is zero. Inserting \(x\) multiplies each weight by
its own critical point, and the sum becomes one. Inserting \(x^3\)
instead gives \(t\). This sphere amplitude reads off the splitting
parameter throughout the family and stays regular at the collision.

This elementary example has an exact converse. Fix a coordinate, a
volume form \(dx\), and a monic critical equation of degree \(d\).
There are \(d\) sphere insertions which recover the entire equation
over the parameter ring. The same measurements determine amplitudes
at every genus. They remain meaningful for multiple critical points
and for nilpotent variations of the potential.

\section{The diagonal and the handle}

Let \(R\) be a commutative \(\mathbb Q\)-algebra, and let
\[
 P(x)=W'(x)=x^d+a_{d-1}x^{d-1}+\cdots+a_0,\qquad W(0)=0,\quad d\ge1 .
\]
The closed algebra \(A=R[x]/(P)\) is free on \(1,x,\ldots,x^{d-1}\).
Write \(\tau(f)\) for the coefficient of \(x^{d-1}\) in the remainder
of \(f\) upon division by \(P\). Equivalently, \(\tau(f)\) is the
coefficient of \(x^{-1}\) in the expansion of \(f(x)/P(x)\) at
infinity. Over \(\mathbb C\) the residue theorem gives
\[
 \tau(f)=\sum_{P(\alpha)=0}
 \Res_{x=\alpha}\frac{f(x)\,dx}{P(x)} .
\]
The formula with remainders defines this trace on every fibre at once.
It fixes the normalization of the protected sphere amplitudes
\cite{vafa}.

To attach a handle one needs the class of the diagonal. In
\(A\otimes_R A\), with coordinates \(x,y\) on its two factors, set
\[
 \mathcal D(x,y)=\frac{P(x)-P(y)}{x-y}.
\]
This divided difference is a polynomial. It satisfies
\((x-y)\mathcal D=0\), and its coefficient of \(x^{d-1}\) is one.
It follows that
\[
 (\tau_x\otimes 1)\bigl(f(x)\mathcal D(x,y)\bigr)=f(y).
\]
Thus \(\mathcal D\) is the inverse of the pairing \(\tau(fh)\).
The Gram matrix gives another way to see perfectness. Its entries
\(\tau(x^{i+j})\) vanish for \(i+j<d-1\) and equal one for
\(i+j=d-1\).

Closing the two legs of the diagonal sets \(y=x\). The resulting
handle element is \(e=P'(x)\). For a product \(f\) of marked
insertions, the genus-\(g\) amplitude is consequently
\begin{equation}
 A_g(f)=\tau\bigl(fP'^{\,g}\bigr),\qquad
 A_{g+1}(f)=\Tr_{A/R}\bigl(m_f m_{P'}^{\,g}\bigr) .
 \label{eq:handle}
\end{equation}
Here \(m_f\) is multiplication by \(f\). For the second equality
choose dual bases for the trace and take the ordinary matrix
trace of multiplication. The contraction joins the same two legs.
Every operation in this calculation commutes with arbitrary changes
of \(R\). The handle has a resolvent over \(R\). Its power series
is \(\sum_g A_g(f)z^g=\tau((1-zm_{P'})^{-1}f)\),
with denominator \(\det(1-zm_{P'})\).
For the convention \(W=x^{D}\), division by \(W'=Dx^{D-1}\)
scales the trace by \(1/D\) and the handle by \(D\) relative to
\(W=x^{D}/D\). The torus amplitude remains \(D-1\).

\section{Reconstructing the marked potential}

Let \(M_r=A_0(x^r)\). The first \(d-1\) moments vanish and
\(M_{d-1}=1\). Multiplication of \(P(x)=0\) by \(x^j\) gives a
recurrence for the remaining moments, or equivalently
\begin{equation}
 \sum_{r\ge0} M_r z^r
  =\frac{z^{d-1}}{1+a_{d-1}z+\cdots+a_0z^d}.
 \label{eq:moments}
\end{equation}
The first \(d\) coefficients following \(M_{d-1}\) determine the
denominator by triangular inversion. Hence the \(d\) sphere
insertions \(x^d,\ldots,x^{2d-1}\) recover every coefficient of \(P\).
Integration recovers \(W\) with the stated constant.

The number of independent measurements is sharp among fixed
polynomial insertions valid over every parameter ring. To see what
an insertion can detect, deform \(P=x^d\) to \(x^d+\epsilon q(x)\)
over \(k[\epsilon]/(\epsilon^2)\), with \(\deg q<d\) and \(k\) a
field of characteristic zero. Expanding the inverse critical equation
at infinity gives
\begin{equation}
 \delta\tau(f)=-\epsilon[x^{2d-1}]\bigl(f(x)q(x)\bigr).
 \label{eq:tangent}
\end{equation}
Each fixed sphere insertion supplies one linear functional on the
\(d\)-dimensional space of \(q\)'s. Fewer than \(d\) have a common
nonzero kernel. There is a stronger conclusion for monomials. If
\(x^k\) with \(d\le k\le 2d-1\) is missing, take \(q=x^{2d-1-k}\).
Among all monomials only \(x^k\) detects this first-order variation.
Arbitrarily high powers cannot replace the missing measurement on
this fibre.

\section{The fibre at a collision}

For a simple root \(\alpha\), contraction onto its idempotent gives
\(A_g(f)_\alpha=f(\alpha)P'(\alpha)^{g-1}\). For a root of
multiplicity \(m>1\), write \(P=(x-\alpha)^m U(x)\), with
\(U(\alpha)\ne0\). The local sphere contribution is the coefficient
of \(z^{m-1}\) in \(f(\alpha+z)/U(\alpha+z)\). The genus-one
contribution is \(m f(\alpha)\) by the local multiplication trace in
\eqref{eq:handle}. At higher genus that trace contains \(P'(\alpha)=0\)
and hence vanishes on the field-valued fibre.

One can follow this vanishing throughout the family
\[
 W_t=\frac{x^{d+1}}{d+1}-tx,\qquad x^d=t,\qquad d\ge2 .
\]
Since \(e=d x^{d-1}\), set \(h=(r+(d-1)g-(d-1))/d\). Reduction of
the insertion \(x^r e^g\) gives, at every genus,
\begin{equation}
 A_g(x^r)=
 \begin{cases}d^g t^h,& h\in\mathbb Z_{\ge0},\\
                0,&\text{otherwise}.\end{cases}
 \qquad
 \sum_{g\ge0} A_g(1)z^g=\frac{d z}{1-d^d t^{d-1}z^d}.
 \label{eq:family}
\end{equation}
For the cubic, the genus-three vacuum amplitude is \(8t\). Taking
\(t=\epsilon\) gives \(8\epsilon\), although the genus-three
amplitude vanishes on the underlying point. For a sum of potentials
in independent coordinates the diagonal and trace are tensor
products. The answer for \(\prod_i x_i^{r_i}\) is the product of
\eqref{eq:family} for the pairs \(d_i,t_i\), including fibres on which
only some of the critical points have coalesced.

The choice of sector matters. Over \(k(t)\) the critical values of
\(W_t\) are \(-dt\alpha/(d+1)\), all nonzero. The polynomial \(W_t\)
and its derivative are coprime. Choose \(u,v\) with \(uW_t+vW_t'=1\).
For a polynomial matrix factorization \(Q^2=W_t\,1\), differentiation
gives \(\{Q,\partial_x Q\}=W_t'\,1\), so
\[
 \bigl\{Q,\tfrac12uQ+v\,\partial_xQ\bigr\}=1 .
\]
The category of finite free factorizations at value zero is
therefore contractible on this generic fibre. The closed theory of
the entire critical scheme still has \(A_1(1)=d\). Comparing it with
matrix factorizations requires selecting the corresponding critical
values and supports. For a fixed isolated critical point the
stabilized boundary, its endomorphism algebra and the identity
defect supply this protected comparison \cite{dyck,dm}.

\begin{thebibliography}{9}\small
\bibitem{vafa} C.~Vafa, \emph{Topological Landau--Ginzburg models},
Modern Phys.\ Lett.\ A \textbf{6} (1991), 337--346.
\href{https://doi.org/10.1142/S0217732391000324}{doi:10.1142/S0217732391000324}.
\bibitem{dyck} T.~Dyckerhoff, \emph{Compact generators in categories of
matrix factorizations}, Duke Math.\ J.\ \textbf{159} (2011), 223--274.
\href{https://arxiv.org/abs/0904.4713}{arXiv:0904.4713}.
\bibitem{dm} T.~Dyckerhoff and D.~Murfet, \emph{The
Kapustin--Li formula revisited}, Adv.\ Math.\ \textbf{231} (2012),
1858--1885. \href{https://arxiv.org/abs/1004.0687}{arXiv:1004.0687}.
\end{thebibliography}
\end{document}
